\chapter{Conclusão}
\label{cap:conclusao}

Este trabalho investigou o que acontece com a autonomia de microsserviços quando eles reutilizam componentes internos. Em vez de partir da percepção de quem desenvolve, como os estudos com praticantes, a análise usou o código e o histórico de oito sistemas Java de código aberto, escolhidos por sorteio a partir de duas bases públicas.

Em relação à RQ1, o compartilhamento é comum, mas não universal. Cerca de metade dos sistemas examinados com três ou mais serviços tem ao menos um componente compartilhado, e, nos sistemas que compartilham, a dependência é ampla: um componente atinge em mediana quatro serviços, e quase nove em cada dez serviços usam algum componente. Em relação à RQ2, há indício de perda de autonomia. Na maioria dos sistemas, os serviços que dependem de um mesmo componente mudam juntos com mais frequência do que os demais, com efeito grande nos dois sistemas em que a diferença foi significativa. Esse indício precisa ser lido com cuidado, porque menos da metade da co-mudança inspecionada tem causa funcional; o resto vem de atualizações em lote que tocam vários serviços sem que um dependa do outro. Em relação à RQ3, o número de serviços que usam um componente e a natureza técnica ou de domínio do seu código não explicaram quanto as mudanças do componente se propagam. As maiores taxas apareceram em módulos pequenos que descrevem a interface de um serviço, enquanto bibliotecas grandes e muito alteradas evoluíram com mais independência. Disso decorre a hipótese de que o compartilhamento compromete a autonomia sobretudo quando o componente é o acordo entre os serviços.

Para quem mantém sistemas de microsserviços, o resultado mais útil é a taxa de propagação. Ela pode ser calculada a partir do histórico de qualquer repositório e mostra, para cada componente, com que frequência mudá-lo exige mudar também os serviços que o usam. Esse número diz mais sobre o custo de um componente compartilhado do que a simples contagem dos seus dependentes.

Os achados têm limites. O corpus tem oito sistemas, vários deles de demonstração, todos em Java e organizados em monorepositórios. A estrutura de módulos atual foi aplicada a todo o histórico, o que reduziu a cobertura em dois sistemas reorganizados. A classificação dos componentes e a inspeção dos \textit{commits} contaram com apoio de modelo de linguagem e revisão da equipe, e a relação entre propagação e tamanho do componente surgiu de uma análise não planejada.

Esses limites indicam os próximos passos. O primeiro é testar, em um corpus maior e com protocolo definido de antemão, a hipótese de que contratos de API propagam mais mudanças do que outros componentes. O segundo é mapear os módulos de cada versão do repositório, e não só da versão atual, para acompanhar sistemas que mudaram de estrutura. O terceiro é estender a análise a outras linguagens, a sistemas distribuídos em vários repositórios e a componentes publicados como pacotes versionados, nos quais a dependência passa a ser mediada por versões.
